\begin{tikzpicture}[font=\small]
  % ---------- the plant ----------
  \begin{scope}
    % cart
    \draw[ssgray, line width=1.0pt] (-2.6,0) -- (2.6,0);
    \foreach \x in {-2.4,-1.8,...,2.4} \draw[sslight, line width=0.6pt] (\x,0) -- (\x-0.25,-0.28);
    \draw[ssgray, line width=1.1pt, fill=sslight!25] (-0.85,0.12) rectangle (0.85,0.85);
    \filldraw[ssgray] (-0.5,0.12) circle (0.12);
    \filldraw[ssgray] (0.5,0.12) circle (0.12);
    % pendulum, tilted
    \draw[curve, line width=2.0pt] (0,0.85) -- (0.72,3.0);
    \filldraw[ssblue] (0.72,3.0) circle (0.20);
    \draw[helper] (0,0.85) -- (0,3.2);
    \draw[ssred, line width=0.8pt] (0,2.1) arc (90:71:1.25);
    \node[ssred, font=\footnotesize, anchor=west] at (0.22,2.25){$\theta$};
    % applied force
    \draw[-{Stealth[length=6pt]}, ssteal, line width=1.3pt] (-2.0,0.5) -- (-0.95,0.5);
    \node[ssteal, font=\footnotesize, anchor=east] at (-2.05,0.5){$f(t)$};

    \node[note, anchor=north, align=center] at (0,-0.55)
      {线性化后 $\ddot\theta-\dfrac{g}{L}\theta=-\dfrac{1}{L}\ddot{s}$};
  \end{scope}

  % ---------- open-loop poles: one in RHP ----------
  \begin{scope}[xshift=6.4cm, yshift=1.3cm]
    \node[anchor=south, font=\small] at (0,2.5){开环极点};
    \draw[ax] (-2.6,0) -- (2.8,0); \node[note, anchor=north west] at (2.6,-0.05){$\sigma$};
    \draw[ax] (0,-2.1) -- (0,2.3); \node[note, anchor=south west] at (0.08,2.15){$j\omega$};
    \foreach \x/\col in {-1.5/ssgray, 1.5/ssred}{
      \draw[\col, line width=1.5pt] (\x-0.16,-0.16) -- (\x+0.16,0.16);
      \draw[\col, line width=1.5pt] (\x-0.16,0.16) -- (\x+0.16,-0.16);
    }
    \node[ssred, font=\footnotesize, anchor=south] at (1.5,0.25){$+\sqrt{g/L}$};
    \node[note, font=\footnotesize, anchor=south] at (-1.5,0.25){$-\sqrt{g/L}$};
    \fill[ssred, opacity=0.10] (0,-2.0) rectangle (2.7,2.0);
    \node[ssred, font=\footnotesize, anchor=north, align=center] at (1.5,-0.35)
      {右半平面极点\\ $\Rightarrow$ 不稳定};
    \node[note, anchor=north, align=center] at (0,-2.4)
      {$H(s)=\dfrac{-1/L}{s^2-g/L}$};
  \end{scope}

  % ---------- proportional vs PD feedback ----------
  \begin{scope}[xshift=13.6cm, yshift=1.3cm]
    \node[anchor=south, font=\small] at (0,2.5){闭环极点的去向};
    \draw[ax] (-2.6,0) -- (2.8,0); \node[note, anchor=north west] at (2.6,-0.05){$\sigma$};
    \draw[ax] (0,-2.1) -- (0,2.3); \node[note, anchor=south west] at (0.08,2.15){$j\omega$};

    % proportional only: poles move onto the jw axis (marginally stable, oscillates)
    \foreach \y in {-1.3,1.3}{
      \draw[ssred, line width=1.4pt] (-0.16,\y-0.16) -- (0.16,\y+0.16);
      \draw[ssred, line width=1.4pt] (-0.16,\y+0.16) -- (0.16,\y-0.16);
    }
    \node[ssred, font=\footnotesize, anchor=west] at (0.28,1.45){纯比例：停在虚轴};

    % PD: poles move into LHP
    \foreach \y in {-1.05,1.05}{
      \draw[ssteal, line width=1.4pt] (-1.5-0.16,\y-0.16) -- (-1.5+0.16,\y+0.16);
      \draw[ssteal, line width=1.4pt] (-1.5-0.16,\y+0.16) -- (-1.5+0.16,\y-0.16);
    }
    \draw[-{Stealth[length=5pt]}, sslight, line width=0.9pt] (-0.3,1.15) -- (-1.2,1.05);
    \node[ssteal, font=\footnotesize, anchor=north] at (-1.5,-1.55){加微分项后\\ 进左半平面};

    \node[note, anchor=north, align=center] at (0,-2.4)
      {比例项只能把极点推到虚轴（等幅摆动）；\\
       必须再引入\textbf{角速度}反馈（微分项）才能提供阻尼};
  \end{scope}

  \node[note, anchor=north, align=center] at (6.4,-3.2)
    {这是本课的收尾例子，也是反馈最有说服力的一处：
     \textbf{反馈能让一个本来不稳定的对象变稳定} ——
     开环右半平面极点被闭环拉进左半平面。\\
     反过来说也成立：反馈用得不当，同样能把稳定的对象变不稳定。同一个分母 $1+GH$ 两面都管。};
\end{tikzpicture}
